% !TeX TS-program = xelatex
% !TEX root = main.tex

% How to:
%


% Load class with all definitions. Do not remove this line.
% Options will be passed to Memoir
\documentclass[final]{multiversum}
%

%%%
% Set those variables

% Authors of the document.
% e.g. Max Mustermann, Erika Musterfrau
\multiauthor{Franca, Rene}

% Date of release.
% e.g. 31.12.2074
\multidate{}

% Number of release, no leading zeros.
% e.g. 15
\multiausgabe{32}

% Losung
% e.g. Die Kuh lief um den Teig.
\multilosung{Wer dient dem Fuchsgott mit der Säge? Der Flex-Geweihte}


% Logo
% Use a different logo. Defaults to Ueberschrift.svg
%\multilogo{Ueberschrift_xmas}

%
%%%

%%%%%%%%%%%%%%%%%%%%%%%%%%%%%%%%% DOCUMENT %%%%%%%%%%%%%%%%%%%%%%%%%%%%%%%%%
\begin{document}

\makemultititle
%

% PUT BODY HERE
\section{Was bisher geschah...}

\subsection{Workshop: ABC}
Lorem Ipsum
\verfasser{Autor1}

\subsection{In der Redaktion}
Lorem Ipsum
\verfasser{Redaktion}

\section{An einem anderen Ort}

\subsection{Eine Prinzessin, etwas Ehre und viele Tote. -- Eine Filmkritik zu \enquote{Das Bündnis}}
\subsubsection{Everybody was konfus fighting}

\enquote{Warum?} - \enquote{Du hast die Prinzessin entehrt!} - \enquote{Was hab ich getan?!}
Diesen Dialog lässt der neue Blockbuster \enquote{Das Bündnis} aus den New Hong Kong Studios seine Protagonisten immer
wieder aufs Neue führen. Das finale Urteil darüber, was die Ehre der Prinzessin so sehr beflecken kann, dass 50 Männer
dafür in den Tod gehen, überlässt der Wuxia-Film dem Publikum.

\enquote{Das Bündnis} beginnt mit einer arrangierten Ehe. Prinzessin Weiter Himmel soll dafür in den Wasserpalast
des Kaisers eskortiert werden. Ihre Eskorte:
Der charismatische Prinz Zhi
(herrlich versnobt gespielt von Kung Fu Star Wei Zhu aka Dennis Wei, dieser demonstriert in \enquote{Das Bündnis} sein
Talent mit dem Wurffächer),
der gemütliche Koch Chao (kumpelhaft-liebenswert dargestellt von Alan \enquote{Iron Fist} Yang,
der, um keinen Fettanzug tragen zu müssen, für die Rolle 100 Kilo zugenommen hat und seinen Wok mit tödlicher Präzision
schwingt)
und der ehrbare Oberaufseher Wang (verkörpert von einem durchtrainierten Danny Lee, dessen Kampfstil mit drei
Kurzschwertern pro Hand nostalgisch an einen jungen Hugh Jackman erinnert).

Einen Sturz aus dem Fenster später ist die Eskorte samt Prinzessin schon unterwegs. Trotz Eile fordert die Prinzessin
direkt eine Pause aufgrund eines Missgeschicks ihres Bruders: Während dieser seine Schwester zur Kutsche trägt,
stolperte er über einen ungünstig platzierten Stein, sodass sowohl Prinz als auch Prinzessin ein unfreiwilliges
Schlammbad nehmen.
In dieser Szene stellt die Prinzessin zum ersten Mal die Kompetenz ihrer Eskorte in Frage, ein Motiv,
was den Rest des Films begleitet.
Erstes Etappenziel ist daher das Badehaus von Boshi, einem alten Freund und Berufskollegen von Koch Chao.
Während Chao und Boshi nostalgische Anekdoten und wohlchoreografierte High Fives austauschen,
liefert sich Wang in der Gaststube ein Blickduell mit einem mysteriösen Unbekannten - und blinzelt als erster.
Foreshadowing oder doch die chronische Bindehautentzündung von Schauspieler Danny Lee?
Bevor die Identität des Unbekannten geklärt werden kann, lassen die Reisegefährten im Badehaus alle Hüllen fallen:
Die Prinzessin, Wang und Chao abends in den heißen Quellen (nicht gemeinsam, sondern sittlich nacheinander),
Prinz Zhi am nächsten Morgen, weil ein mongolischer Schurke ihm nachts Hab, Gut und Hose entwendet hat.

Aber \enquote{Das Bündnis} ist nicht nur visuell ansprechend, sondern verbindet Action und Humor:
Ninjas attackieren in einem Bambuswald die Kutsche mit der Prinzessin.
Prinz Zhi zieht die Wurffächer, Koch Chao Wok und Kochlöffel; nur Oberaufseher Wang flitzt ehrenhaft blind
mit seinen sechs Schwertern an der Kampfsituation vorbei (Schauspieler Lee trotz Bindehautentzündung professionell wie immer.)
Später stemmt Chao die Kutsche wegen eines Radbruchs mittels einer Bambusstange hoch.
Leider spart \enquote{Das Bündnis} hier am Stunt-Budget, sodass der flexible Bambus weder den fülligen Koch
noch Prinz Zhi samt Pferd in einen Orbit katapultiert.
Kritik auf hohem Niveau: Die Bambusszenen sind ein Highlight des Films.

Das Finale des Films führt die Reisegefährten an den Hafen von Pulin.
Dort wartet ein Schiff auf die Prinzessin, aber auch der Unbekannte aus Boshis Gaststube,
diesmal mit zwei Freunden.
Fliegender Adler, so sein Name, gibt sich als der heimliche Geliebte der Prinzessin zu erkennen,
der die arrangierte Ehe verhindern und mit ihr durchbrennen will,
doch ist seine Familie mit der Sippschaft der Prinzessin verfeindet.
Und als wäre die Tragödie noch nicht vertrackt genug, taucht auch noch der ältere Bruder von Fliegender Adler auf,
wenig begeistert über die Affäre und den Betrug seines Bruders.
Doch bevor das Publikum verstehen kann, wer nun die Bösen und die Guten sind,
tritt der ehrbare Wang vor und fordert Fliegender Adler zum Duell.
Endlich spielt das musikalische Thema, das Fans inzwischen mit Danny Lees Rollen assoziieren:
der ikonische Text \enquote{Let's get down to business/ to defeat the huns}
wird hier interpretiert von einem Chor tibetanischer Schweigemönche,
sodass der Fokus eher auf der instrumentalen Wucht des Musikstücks liegt.
Eine Hommage mit feiner Ironie an die hosenstehlenden Hunnen im Badehaus - diese mögen zwar Gauner sein,
aber die wahren Schurken des Konflikts sind sie nicht.

Doch wer sind hier die Bösen und die Guten? Schnell sind in den Kampf am Hafen
nicht nur Wang und Fliegender Adler verwickelt, sondern auch der Bruder, Chao, Prinz Zhi und etwa 50 Statisten.
Hier offenbart \enquote{Das Bündnis} erneut sein humoristisches Talent.
Auf dem schmalen Grat zwischen Tragödie und Slapstick verwirbeln Stunts, Unfälle und Kollateralschaden zu einer
choreografierten Katastrophe.
Faustschläge verfehlen den Gegner und treffen den Verbündeten, Kleidung wird in Brand gesteckt,
Kämpfer ertrinken im eisigen Flusswasser, Waffen enden im eigenen Bein.
Ist es Versehen oder ist es Absicht -- dem Publikum bleibt keine Zeit, bei diesen waghalsigen Schnitten,
Close-Ups und Überblenden darüber nachzudenken.
Wang zielt auf den Bruder und trifft Fliegender Adler, Chao schleudert ein Küchenutensil auf Wang,
Fliegender Adler attackiert Wang und verletzt seinen Bruder,
und immer wieder erklingt nach einem schweren Treffer der empörte Dialog:
\enquote{Warum?} - \enquote{Du hast die Prinzessin entehrt!} - \enquote{Was hab ich getan?!}
Längst hat das Publikum gemeinsam mit den Protagonisten die Übersicht verloren,
wer denn nun wie und warum die Prinzessin entehrt hat.
Durch dieses Hin und Her scheint die Bewachung der Prinzessin zwischenzeitlich völlig vergessen,
sodass die beiden Freunde von Fliegender Adler diese Aufgabe übernehmen müssen,
bei der auch sie schlussendlich ihr Leben lassen.
Als Finale wirft Wang sein Schwert in die Luft,
wird selbst davon getroffen und führt die ikonischen drei Zeilen als Selbstgespräch.
Während dieser ergreifenden Einstellung wird Wangs Sixpack ein letztes Mal in Szene gesetzt,
was die Fans von Danny Lee freuen wird.
Die Prinzessin selbst konstatiert -- mit Recht -- wiederholt die Inkompetenz aller Beteiligten.

Als der Kampf vorbei ist, färbt sich das Wasser des Flusses rot.
Wie die Reise für Wang und Chao mit einem fulminanten Fenstersturz begann,
so endet sie für Wang und Fliegender Adler mit einem Salto Mortale vom Dach.
Ein Schiff mit der Prinzessin, Prinz Zhi und Chao segelt dem unausweichlichen Schicksal entgegen,
das titelgebende Bündnis -- die arrangierte Ehe -- wird endlich vollzogen,
es ist das bittersüße Ende einer rasanten Martial Arts Tragikomödie.
Und man erinnert sich an die letzten Worte des sterbenden Wang zu seinem Freund Chao:
\enquote{We must be swift as the coursing river / With all the force of a great typhoon / Be a man!}
und unterdrückt die Tränen nur, weil er in einer Post-Credits-Szene Chao sein geheimes Knödelrezept preisgibt.
\verfasser[New Hong Kong Story]{Franca \& Rene}

\subsection{Titel 2}
Lorem Ipsum, Text endet an Zeilenende.
\Verfasser[System]{Autor2}


\begin{termine}
% Put dates here:
  \item Monatstreffen: 16.09.25, 19 Uhr
  \item Stadtführung durch Köln (Sagen \& Legenden): 30.08.25 (Uhrzeit folgt auf Webseite)
  \item SchwefeldrAachen Con: 13.09.25 - 14.09.25
  \item Freizeit Freusburg: 10.01.26 - 14.01.26
\end{termine}
\impressum

\end{document}
%%%%%%%%%%%%%%%%%%%%%%%%%%%%%%%%% END DOCUMENT %%%%%%%%%%%%%%%%%%%%%%%%%%%%%%%%%
